\documentclass[12pt]{article}
\usepackage[utf8]{inputenc}
\usepackage[T1]{fontenc}
\usepackage{hyperref}
\title{Parthian Language Resources for K-12 Curriculum Development}
\author{Compiled by Claude Sonnet 5.5 (claude-sonnet-5-5) for Andr\'as Kornai}
\date{}
\begin{document}
\maketitle

\section{Introduction}
Parthian (ISO 639-3 xpr; Glottocode part1239) is an extinct Middle Iranian language of the Indo-European family. It has no native speakers. The search found only scholarly resources for it. It did not find any school curriculum, primer or revitalization programme. Any K-12 heritage curriculum would therefore have to be built from academic sources.

\section{Primary Educational Resources}
No primers, textbooks or school programmes for Parthian turned up in the search. The nearest things to teaching material are the scholarly readers and grammars listed below. They would need adapting for school use.

\section{Linguistic and Reference Materials}
\begin{itemize}
\item Desmond Durkin-Meisterernst, \emph{Grammatik des Westmitteliranischen (Parthisch und Mittelpersisch)}, Grammatica Iranica 1, Austrian Academy of Sciences, 2014. The grammar describes Parthian and Middle Persian together, covering scripts, phonology, morphology and syntax. The publisher page is \url{https://verlag.oeaw.ac.at/en/product/grammatik-des-westmitteliranischen/30312}. A further description is at \url{https://biblioiranica.info/grammatik-des-westmitteliranischen-parthisch-und-mittelpersisch}.
\item A Cambridge article on a Parthian sound change, published in the Bulletin of SOAS: \url{https://resolve.cambridge.org/core/journals/bulletin-of-the-school-of-oriental-and-african-studies/article/footnotes-on-a-parthian-sound-change1/45D508D8C8AD5CF92B2C6AF4F53840B1}.
\item A preliminary report on the classification of Middle Iranian languages, from Masaryk University: \url{https://digilib.phil.muni.cz/en/node/27020}.
\end{itemize}
The search results also mentioned the following printed works, but gave no URLs for them. These are Durkin-Meisterernst's \emph{Dictionary of Manichaean Middle Persian and Parthian} (Brepols, 2004), Mary Boyce's \emph{A Reader in Manichaean Middle Persian and Parthian} (1975), and Boyce's \emph{A Word-List of Manichaean Middle Persian and Parthian} (1977).

\section{Religious and Community Resources}
The sources mention Manichaean Parthian texts, including hymn cycles. No community organization or revitalization effort for Parthian was found.

\section{Digital and Audio Resources}
No audio, video or software resources were found in the search results.

\section{Curriculum Development Context}
Parthian is a classical, extinct language that is documented only in scholarly literature. A school course would most plausibly be a heritage and history unit, with Parthian taught through short Manichaean texts and an introduction to the script. The published grammar and readers would serve as the teacher's reference.
\end{document}